\documentclass[12pt]{article}
\usepackage[utf8]{inputenc}
\usepackage[T1]{fontenc}
\usepackage[margin=1in]{geometry}
\usepackage{setspace}
\usepackage{xcolor}
\usepackage[colorlinks=true, urlcolor=blue!60!black, linkcolor=blue!60!black]{hyperref}
\setlength{\parindent}{2em}
\setlength{\parskip}{1.25em}
\renewcommand{\baselinestretch}{1.0}
\title{\textit{Kelo v. City of New London}, 545 U.S. 469 (2005)}
\author{Case Brief}
\date{}
\begin{document}
\maketitle

\section{Case Metadata}
\begin{itemize}
  \item \textbf{Case Name and Citation:} \textit{Kelo v.\ City of New London}, 545 U.S.\ 469 (2005) (No.\ 04-108)
  \item \textbf{Argued / Decided:} Argued February 22, 2005; decided June 23, 2005
  \item \textbf{Vote:} 5--4
  \item \textbf{Author:} Justice John Paul Stevens (joined by Justices Kennedy, Souter, Ginsburg, and Breyer)
  \item \textbf{Concurrences / Dissents:} Concurrence by Justice Anthony M.\ Kennedy; dissent by Justice Sandra Day O'Connor (joined by Chief Justice Rehnquist and Justices Scalia and Thomas); separate dissent by Justice Clarence Thomas
  \item \textbf{Chapter / Topic:} Chapter 6: Substantive Due Process --- Takings \& Property Rights (Fifth Amendment Public Use Clause)
  \item \textbf{Primary Source URLs:} \href{https://supreme.justia.com/cases/federal/us/545/469/}{Justia U.S.\ Supreme Court Center} $|$ \href{https://www.law.cornell.edu/supct/html/04-108.ZO.html}{Cornell LII}
\end{itemize}

\section{Facts and Procedural History}
Following decades of economic decline---including a 1990 state designation as a ``distressed municipality'' and the 1996 federal closure of the Naval Undersea Warfare Center, which eliminated over 1,500 jobs---New London, Connecticut, had an unemployment rate nearly double the state average and its lowest population since 1920. To revitalize the 90-acre Fort Trumbull peninsula, city officials reactivated the New London Development Corporation (NLDC), a private nonprofit entity, while Connecticut authorized \$5.35 million in planning bonds and \$10 million for a new Fort Trumbull State Park. In February 1998, pharmaceutical giant Pfizer Inc.\ announced plans to build a \$300 million research facility immediately adjacent to Fort Trumbull. Seeking to capitalize on Pfizer's arrival, the NLDC drafted a seven-parcel integrated redevelopment plan featuring a waterfront conference hotel, restaurants, a riverwalk, a renovated marina, 80 new residences, space for a U.S.\ Coast Guard Museum, at least 90,000 square feet of research-and-development office space (Parcel~3), and a 2.4-acre ``park support'' or marina-support tract (Parcel~4A).

In January 2000, the city council approved the plan under Connecticut's municipal development statute (Conn.\ Gen.\ Stat.\ \S~8-186 \textit{et seq.}) and authorized the NLDC to acquire the remaining parcels by eminent domain. After purchasing most of the area's 115 private properties from willing sellers, the NLDC initiated condemnation proceedings against nine unwilling owners of 15 well-maintained, unblighted properties (4 in Parcel~3 and 11 in Parcel~4A). Petitioners included Susette Kelo, who had extensively renovated her waterfront home since 1997, and Wilhelmina Dery, who was born in her Fort Trumbull home in 1918 and had lived there her entire life.

In December 2000, the homeowners sued in New London Superior Court, alleging that condemning nonblighted homes for private economic development violated the Fifth Amendment's Public Use Clause (incorporated via the Fourteenth Amendment). After a seven-day bench trial, the Superior Court enjoined the takings on Parcel~4A as too vague, but upheld the takings on Parcel~3. On cross-appeals, the Connecticut Supreme Court (4--3) upheld all the condemnations under \textit{Berman v.\ Parker} and \textit{Hawaii Housing Authority v.\ Midkiff}, holding that economic development constituted a valid public use under both the state and federal constitutions (268 Conn.\ 1, 843 A.2d 500 (2004)). The U.S.\ Supreme Court granted certiorari.

\section{Issue(s) Presented}
Does a city's condemnation of unblighted private property for transfer to private developers as part of a comprehensive economic-development plan satisfy the ``public use'' requirement of the Fifth Amendment's Takings Clause?

\section{Holding}
Yes. Condemning private property for transfer to private parties pursuant to a carefully formulated, comprehensive economic-development plan serves a valid ``public purpose'' and therefore satisfies the Fifth Amendment's ``public use'' requirement under deferential judicial review, even where the condemned properties are not blighted and will not be opened to general public use. States remain free to impose stricter public-use limits under state constitutional or statutory law. Judgment affirmed.

\section{Reasoning}
Writing for a five-Justice majority, Justice Stevens situated New London's plan squarely within a century of deferential Public Use Clause precedent:

\begin{itemize}
  \item \textbf{Rejection of Literal ``Use by the Public'':} Stevens began with two settled poles: government may not take property from $A$ solely to transfer it to $B$ for private benefit (or under the ``mere pretext'' of a public purpose), yet it may transfer property to a private common carrier open to the public. Between those poles, the Court had ``long ago rejected any literal requirement that condemned property be put into use for the general public,'' embracing instead the broader concept of ``public purpose'' (\textit{Fallbrook Irrigation Dist.\ v.\ Bradley}, 164 U.S.\ 112 (1896); \textit{Strickley v.\ Highland Boy Gold Mining Co.}, 200 U.S.\ 527 (1906)).
  \item \textbf{Deference Under \textit{Berman} and \textit{Midkiff}:} Relying on \href{https://supreme.justia.com/cases/federal/us/348/26/}{\textit{Berman v.\ Parker}}, 348 U.S.\ 26 (1954) (upholding condemnation of an unblighted department store within an integrated D.C.\ urban-renewal plan), \href{https://supreme.justia.com/cases/federal/us/467/229/}{\textit{Hawaii Housing Authority v.\ Midkiff}}, 467 U.S.\ 229 (1984) (upholding land-redistribution transfers from lessors to lessees to break a land oligopoly), and \textit{Ruckelshaus v.\ Monsanto Co.}, 467 U.S.\ 986 (1984), Stevens emphasized that the Court evaluates comprehensive redevelopment plans as a whole rather than lot by lot, focusing on the taking's \textit{purpose} rather than its private-transfer \textit{mechanics}.
  \item \textbf{Economic Development as a Traditional Public Purpose:} Because promoting economic rejuvenation, job creation, and tax-base growth is a traditional governmental function, Stevens rejected petitioners' proposed bright-line ban on economic-development takings. Because New London acted pursuant to a thoroughly deliberated, state-authorized plan where most private beneficiaries were not yet identified, there was no evidence of illegitimate private favoritism.
  \item \textbf{Rejection of ``Reasonable Certainty'' Review and Federalism Coda:} Declining to require proof to a ``reasonable certainty'' that projected economic gains would materialize, Stevens warned that federal courts must not second-guess socioeconomic predictions or project boundaries. He closed by stressing that the Fifth Amendment sets only a federal floor: ``nothing in our opinion precludes any State from placing further restrictions on its exercise of the takings power.''
\end{itemize}

\section{Concurring and Dissenting Opinions}
\begin{itemize}
  \item \textbf{Justice Kennedy (concurring):} Joining the majority in full, Kennedy emphasized that rational-basis review under the Public Use Clause still bars takings ``intended to favor a particular private party, with only incidental or pretextual public benefits'' (analogizing to equal-protection animus cases like \textit{Cleburne} and \textit{Moreno}). Reviewing the trial record, Kennedy found no pretext here because New London faced genuine economic distress, committed public funds before most private beneficiaries were known, and followed elaborate planning procedures. He left open whether ``a more stringent standard of review'' might apply to more suspicious one-to-one private transfers.
  \item \textbf{Justice O'Connor (dissenting, joined by Chief Justice Rehnquist and Justices Scalia and Thomas):} O'Connor warned that equating ``public use'' with incidental economic benefits ``effectively delete[s] the words `for public use' from the Takings Clause.'' Distinguishing \textit{Berman} and \textit{Midkiff} as cases where the \textit{precondemnation use} itself inflicted affirmative social harm (slum blight or land oligopoly), she argued that New London's upgrade rationale leaves every home, farm, or small business vulnerable to transfer to wealthier, politically connected developers, while Kennedy's pretext inquiry would be failed only by a ``stupid staffer.''
  \item \textbf{Justice Thomas (dissenting):} Joining O'Connor's dissent, Thomas urged overruling \textit{Berman} and \textit{Midkiff} altogether to restore the Public Use Clause's original meaning: government may condemn property only if the government owns it or the public has a legal right to use it. Pointing to the history of mid-century urban renewal (``Negro removal''), Thomas argued that deferential ``public purpose'' review perversely exposes politically powerless and minority communities to disproportionate displacement.
\end{itemize}

\section{Key Quotations}
\begin{enumerate}
  \item ``Promoting economic development is a traditional and long accepted function of government. There is, moreover, no principled way of distinguishing economic development from the other public purposes that we have recognized. \dots\ We emphasize that nothing in our opinion precludes any State from placing further restrictions on its exercise of the takings power.'' --- Justice Stevens, 545 U.S.\ at 484, 489 (\href{https://supreme.justia.com/cases/federal/us/545/469/}{Justia}).
  \item ``A court applying rational-basis review under the Public Use Clause should strike down a taking that, by a clear showing, is intended to favor a particular private party, with only incidental or pretextual public benefits.'' --- Justice Kennedy (concurring), 545 U.S.\ at 491 (\href{https://supreme.justia.com/cases/federal/us/545/469/}{Justia}).
  \item ``The specter of condemnation hangs over all property. Nothing is to prevent the State from replacing any Motel 6 with a Ritz-Carlton, any home with a shopping mall, or any farm with a factory. \dots\ The beneficiaries are likely to be those citizens with disproportionate influence and power in the political process, including large corporations and development firms.'' --- Justice O'Connor (dissenting), 545 U.S.\ at 503, 505 (\href{https://supreme.justia.com/cases/federal/us/545/469/}{Justia}).
  \item ``The most natural reading of the Clause is that it allows the government to take property only if the government owns, or the public has a legal right to use, the property. \dots\ Though citizens are safe from the government in their homes, the homes themselves are not.'' --- Justice Thomas (dissenting), 545 U.S.\ at 508, 518 (\href{https://supreme.justia.com/cases/federal/us/545/469/}{Justia}).
\end{enumerate}

\section{Doctrinal Significance and Modern Legacy (Through 2026)}
Few modern Supreme Court decisions have provoked a swifter or broader political backlash than \textit{Kelo}. Ironically, New London's Fort Trumbull redevelopment collapsed in practice: developers failed to secure financing, Pfizer closed its New London research center in 2009 as its tax breaks expired, and the condemned neighborhood sat largely vacant for nearly two decades despite roughly \$78 million in public spending (\href{https://ij.org/case/kelo/}{Institute for Justice}).

\begin{itemize}
  \item \textbf{Post-\textit{Kelo} State Reforms:} Taking up Stevens's federalism invitation, \href{https://ij.org/issues/private-property/eminent-domain/}{47 states} enacted statutory, constitutional, or judicial reforms curbing economic-development condemnations---including 12 state constitutional amendments (such as \href{https://www.rstreet.org/commentary/texas-and-eminent-domain-law-better-than-elsewhere-but-still-room-for-improvement/}{Texas Proposition 11} in 2009, amending Tex.\ Const.\ art.\ I, \S~17(b) by an 89\% vote to bar takings transferred to private entities primarily for economic development or tax revenue). Scholars note, however, that many state reforms remain porous due to expansive statutory definitions of ``blight.''
  \item \textbf{Continuing Supreme Court Calls to Reconsider \textit{Kelo}:} Through 2026, \textit{Kelo} remains formally binding federal law, even as the Roberts Court has steadily fortified property rights in adjacent Takings Clause decisions (e.g., \textit{Knick v.\ Township of Scott}, 588 U.S.\ 180 (2019); \textit{Cedar Point Nursery v.\ Hassid}, 594 U.S.\ 139 (2021)). In \href{https://www.supremecourt.gov/opinions/20pdf/20-1214_7lh8.pdf}{\textit{Eychaner v.\ City of Chicago}}, 594 U.S.\ \_\_\_ (2021) (involving Chicago's condemnation of unblighted land for transfer to a chocolate manufacturer), Justice Thomas---joined by Justice Gorsuch---dissented from denial of certiorari to urge reconsidering \textit{Kelo} and clarifying its pretext rule, with Justice Kavanaugh also noting he would have granted review. Although the Court subsequently denied certiorari in \href{https://ij.org/case/bowers-v-oneida-county-industrial-development-agency/}{\textit{Bowers v.\ Oneida County Industrial Development Agency}} (March 2025), petitions urging the Court to overrule or tighten \textit{Kelo} continue to reach the docket in 2026.
\end{itemize}

\section{Pedagogical Critiques and Exam Hypotheticals}
\textbf{Pedagogical Critiques:}
\begin{itemize}
  \item \textbf{Just Compensation Does Not Make Owners Whole:} Stevens assumed that ``just compensation'' mitigates condemnation hardship (545 U.S.\ at 489 n.\ 21), yet fair-market-value payouts omit subjective attachment, relocation disruption, and assembly surplus---creating a windfall for private transferees at displaced homeowners' expense.
  \item \textbf{The Unworkability of Kennedy's Pretext Test:} Lower courts after \textit{Kelo} have struggled to apply Kennedy's fact-intensive pretext inquiry, because private profit and public tax-base gains are inextricably intertwined in nearly every economic-development project.
\end{itemize}

\textbf{Exam Hypothetical:} Without any prior city-wide redevelopment study, a prosperous suburb receives a proposal from MegaTech Corp.\ to build a corporate campus if the city acquires an adjacent, profitable family-owned nursery whose owner refused MegaTech's private purchase offer. Two weeks later, the city council passes a one-parcel ``economic revitalization'' resolution condemning the nursery for immediate transfer to MegaTech, citing projected property-tax increases. The nursery owner sues under the Fifth Amendment's Public Use Clause (assume no stricter state-law bar).
\begin{itemize}
  \item \textbf{City's Best Argument (Broad Reading of \textit{Kelo}, \textit{Berman}, and \textit{Midkiff}):} Under \textit{Kelo}, generating jobs and tax revenue is a legitimate public purpose, courts owe extreme deference to legislative judgments, and neither blight nor public access is constitutionally required. As long as the council rationally believes MegaTech's campus will boost municipal revenue, federal courts should not second-guess the condemnation's wisdom or mechanics.
  \item \textbf{Owner's Best Argument (Opposite-Side Contrast via Stevens's Caveat and Kennedy's Concurrence):} Even accepting \textit{Kelo} as binding law, this condemnation crosses the line both Stevens and Kennedy drew. Unlike New London's multi-year, 90-acre integrated plan in a distressed city where most private tenants were initially unknown, this is an isolated, one-to-one transfer ($A$ to $B$) initiated \textit{after} a specific private beneficiary failed to buy the land on the open market, with no comprehensive plan and rushed procedures. Under Stevens's ``mere pretext'' prohibition (545 U.S.\ at 478, 487) and Kennedy's meaningful rational-basis concurrence (545 U.S.\ at 491--493), a court should invalidate the taking as an unconstitutional private-benefit transfer---and preserve an alternative argument to overrule \textit{Kelo} under Justice Thomas's originalist ``actual public use'' standard.
\end{itemize}

\end{document}
